\begin{examplebox}
	\textbf{\textcolor{CExample}{例 1（基础）}}
	\textbf{\textcolor{CExample}{题目：}}
	在 $\triangle ABC$ 中，$a=2,\ b=3,\ A=30^\circ$，判断三角形解的个数。
	\textbf{\textcolor{CExample}{解题步骤：}}
	\begin{description}[style=nextline, leftmargin=2.8em, labelsep=0.5em, itemsep=0.45em]
		\item[\textbf{第一步（计算高 h）：}]
			计算 $h = b\sin A = 3 \times \sin30^\circ = 3 \times \frac12 = 1.5$。
			\par\textbf{理由：} $h = b\sin A$ 是判据的必要参数。
		\item[\textbf{第二步（比较 a 与 h）：}]
			$a=2 > h=1.5$，排除 $a<h$ 和 $a=h$ 的情形。
			\par\textbf{理由：} 当 $a>h$ 时可能有一解或两解。
		\item[\textbf{第三步（比较 a 与 b 并判断）：}]
			因为 $a=2 < b=3$ 且 $A=30^\circ<90^\circ$，所以 $h<a<b$，故有两解。
			\par\textbf{理由：} 锐角条件下 $h<a<b$ 对应两解。
	\end{description}
	\textbf{\textcolor{CExample}{关键结论：}}
	三角形有两解。验证：$\sin B = \frac{b\sin A}{a}=0.75$，$B\approx48.6^\circ$ 或 $131.4^\circ$，均满足内角和。

	\medskip
	\textbf{\textcolor{CExample}{例 2（稍变形）}}
	\textbf{\textcolor{CExample}{题目：}}
	在 $\triangle ABC$ 中，$a=2,\ b=4,\ A=30^\circ$，判断解的个数。
	\textbf{\textcolor{CExample}{解题步骤：}}
	\begin{description}[style=nextline, leftmargin=2.8em, labelsep=0.5em, itemsep=0.45em]
		\item[\textbf{第一步（计算高 h）：}]
			$h = 4 \times \sin30^\circ = 4 \times \frac12 = 2$。
			\par\textbf{理由：} 计算判据中的关键量。
		\item[\textbf{第二步（比较 a 与 h）：}]
			$a=2 = h$，满足 $a=h$。
			\par\textbf{理由：} $a=h$ 时 $\sin B = 1$，$B=90^\circ$。
		\item[\textbf{第三步（下结论）：}]
			根据判据，$A<90^\circ$ 且 $a=h$，故三角形有一解。
			\par\textbf{理由：} 判据中 $a=h$ 对应一解。
	\end{description}
	\textbf{\textcolor{CExample}{关键结论：}}
	一解。由正弦定理得 $\sin B = \frac{b\sin A}{a}=1$，$B=90^\circ$，$C=60^\circ$，三角形唯一。

	\medskip
	\textbf{\textcolor{CExample}{例 3（综合应用）}}
	\textbf{\textcolor{CExample}{题目：}}
	在 $\triangle ABC$ 中，$a=1,\ b=\sqrt{3},\ A=30^\circ$，判断解的个数并求边 $c$ 的长。
	\textbf{\textcolor{CExample}{解题步骤：}}
	\begin{description}[style=nextline, leftmargin=2.8em, labelsep=0.5em, itemsep=0.45em]
		\item[\textbf{第一步（判据判断）：}]
			计算 $h = \sqrt{3} \cdot \frac12 = \frac{\sqrt{3}}{2} \approx 0.866$。因为 $a=1 > h$，$a=1 < \sqrt{3}$，$A=30^\circ<90^\circ$，所以满足 $h<a<b$，有两解。
			\par\textbf{理由：} 直接应用判据。
		\item[\textbf{第二步（求角 B）：}]
			由正弦定理 $\sin B = \frac{b\sin A}{a} = \frac{\sqrt{3}\cdot\frac12}{1} = \frac{\sqrt{3}}{2}$。因为 $A<90^\circ$ 且 $a<b$，$B$ 可取 $60^\circ$ 或 $120^\circ$。
			\par\textbf{理由：} $\sin B$ 已知，结合条件得两角。
		\item[\textbf{第三步（求边 c）：}]
			利用正弦定理 $c = \frac{a\sin C}{\sin A}$。当 $B=60^\circ$ 时 $C=90^\circ$，$c_1 = \frac{1\cdot\sin90^\circ}{\sin30^\circ}=2$；当 $B=120^\circ$ 时 $C=30^\circ$，$c_2 = \frac{1\cdot\sin30^\circ}{\sin30^\circ}=1$。
			\par\textbf{理由：} $C=180^\circ-A-B$，再用正弦定理。
	\end{description}
	\textbf{\textcolor{CExample}{关键结论：}}
	三角形有两解，$c=2$ 或 $c=1$。
\end{examplebox}
